\subsection{INNER AUTOMORPHISMS}

Let \(G\) be a group. The set \(\operatorname{Aut}(G)\) of automorphisms of the group \(G\) is a subgroup of \(\mathfrak{S}_{\mathbb{G}}\) (§ 4, no. 1, Example 2).

PROPOSITION 3. Let G be a group. For every element x of G, the mapping \(\operatorname{Int}(x): y \mapsto xyx^{-1}\) of G into itself is an automorphism of G. The mapping \(\operatorname{Int}: x \mapsto \operatorname{Int}(x)\) of G into \(\operatorname{Aut}(G)\) is a group homomorphism, whose kernel is the centre of G and whose image is a normal subgroup of \(\operatorname{Aut}(G)\) .

If x, y and z are elements of G, then \((xyx^{-1})(xzx^{-1}) = xyzx^{-1}\) and hence \(\operatorname{Int}(x)\) is an endomorphism of G. For x and y elements of G,

\[
\operatorname{Int} (x) \circ \operatorname{Int} (y) = \operatorname{Int} (x y):
\]

for all \(z \in G\) , \(x(yzy^{-1})x^{-1} = (xy)z(xy)^{-1}\) . On the other hand, \(\operatorname{Int}(e)\) is the identity mapping of G. The mapping Int is therefore a monoid homomorphism from G to the monoid \(\operatorname{End}(G)\) of endomorphisms of the group G. As the elements of G are invertible, the mapping Int takes its values in the set \(\operatorname{Aut}(G)\) of invertible elements of \(\operatorname{End}(G)\) (§ 2, no. 3). Now \(xyx^{-1} = y\) if and only if x and y commute and hence \(\operatorname{Int}(x)\) is the identity mapping of G if and only if x is a central element. Finally, let \(\alpha\) be an automorphism of G and let \(x \in G\) ; then

\[
\operatorname{Int} (\alpha (x)) = \alpha \circ \operatorname{Int} (x) \circ \alpha^ {- 1}.\tag{2}
\]

For \(y \in G\) ,

\[
\alpha (x). y. \alpha (x) ^ {- 1} = \alpha (x). \alpha (\alpha^ {- 1} (y)). \alpha (x) ^ {- 1} = \alpha (x. \alpha^ {- 1} (y). x ^ {- 1}).
\]

Hence \(\alpha\) . Int(G). \(\alpha^{-1} \subset \operatorname{Int}(G)\) .

DEFINITION 4. Let G be a group and \(x \in \mathbf{G}\) . The automorphism \(y \mapsto xyx^{-1}\) is called the inner automorphism of G defined by \(x\) and is denoted by Int \(x\) .

For \(x, y \in G\) , we also write \(x^y = y^{-1}xy = (\operatorname{Int} y^{-1})(x)\) .

A subgroup of G is normal if and only if it is stable under all inner automorphisms of G (§ 4, no. 4, Definition 5). A subgroup of G is called characteristic if it stable under all automorphisms of G. The centre of a group G is a characteristic subgroup (formula (2)).

The centre of a group G is not necessarily stable under all endomorphisms of G (Exercise 22). In particular, the centre of a group with operators is not necessarily a stable subgroup.

PROPOSITION 4. Let G be a group, H a characteristic (resp. normal) subgroup of G and K a characteristic subgroup of H. Then K is a characteristic (resp. normal) subgroup of G.

The restriction to H of an automorphism (resp. inner automorphism) of G is an automorphism of H and therefore leaves K invariant.

Let G be a group, A ⊂ G and \(b \in G\) . b is said to normalize A if \(bAb^{-1} = A\) ; b is said to centralize A if, for all \(a \in A\) , \(bab^{-1} = a\) . Let A and B be subsets of G; B is said to normalize (resp. centralize) A if every element of B normalizes (resp. centralizes) A.

The set of \(g \in G\) which normalize (resp. centralize) A is called the normalizer (resp. centralizer) of A (cf.~§1, no. 5, Definition 9); it is often denoted by \(N_{G}(A)\) or simply \(N(A)\) (resp. \(C_{G}(A)\) or \(C(A)\) ). It is a subgroup of G. When A is a subgroup of G, \(N_{G}(A)\) may be characterized as the largest subgroup of G which contains A and in which A is normal.

Remarks. (1) The normalizer (resp. centralizer) of A is the strict stabilizer (resp. fixer) of A when G operates on itself by inner automorphisms. In particular the centralizer is a normal subgroup of the normalizer.

\begin{enumerate}
\def\labelenumi{(\arabic{enumi})}
\setcounter{enumi}{1}
\tightlist
\item
  The set of elements \(b \in \mathbf{G}\) such that \(b\mathbf{A}b^{-1} \subset \mathbf{A}\) is a submonoid of \(\mathbf{G}\) . Even when \(\mathbf{A}\) is a subgroup of \(\mathbf{G}\) , this set is not necessarily a subgroup of \(\mathbf{G}\) (Exercise 27).
\end{enumerate}

